
\begin{proposition}[The reduced horizontal case is a minimal elliptic ruled double-curve fibre]\label{cdp:st:C.2.IRR.1}
Let $f:Y\to C_0$ be a proper holomorphic map from a smooth compact
complex threefold with $H^2(Y,\mathbb Q)=0$, with connected fibres
and smooth general fibre a complex two-torus. Let $mS$ be a full
fibre with irreducible nonnormal reduced support. Let
$h=\eta\pi:Z\to T\to S$ be the minimal resolution of its
normalization, and suppose that $p:Z\to B$ is an irrational
ruling. Write the effective divisor of \ref{cdp:st:A.5.3} as
$D_Z=H+V$, where $V$ is vertical and the horizontal part $H$
is reduced, of total degree two over $B$.

Then $B$ is elliptic, $T=Z$ is a smooth minimal ruled surface
over $B$, and $D_Z=H$ is a disjoint union of smooth elliptic
curves. The reduced conductor image is a disjoint union of
smooth elliptic curves. Its inverse image $H$ maps onto it
by a finite etale covering of total degree two over each
component. The support $S$ has only ordinary double-curve
normal crossings and
\[
                         e(S)=0.                 \tag{C.IR.1}\label{cdp:eq:C.IR.1}
\]
No assumption on the initial normalization singularities or
on the ambient Picard pullback is made.
\end{proposition}
\begin{proof}
All divisor intersections in this proof use rational Chern
classes. The actual canonical comparison in \ref{cdp:st:A.5.3} gives
\[
 K_Z+D_Z=h^*\omega_S,\qquad c_1(K_Z+D_Z)=0.
                                                     \tag{C.IR.2}\label{cdp:eq:C.IR.2}
\]
For an irreducible component $H_i$ of the reduced horizontal
part, adjunction therefore gives
\[
 2p_a(H_i)-2+(D_Z-H_i)\cdot H_i=0.                 \tag{C.IR.3}\label{cdp:eq:C.IR.3}
\]
The second term is nonnegative. The normalization of $H_i$
maps finitely onto the positive-genus curve $B$; Riemann--Hurwitz
makes its genus at least one. Equality in \eqref{cdp:eq:C.IR.3} thus makes
$H_i$ a smooth elliptic curve, makes $B$ elliptic, and makes
all horizontal components disjoint from each other and from
$V$.

We give the required comparison under the ruling blowdowns.
For one $\sigma:W\to W'$ with exceptional curve $E$, put
$D'=\sigma_*D_W$, let $a$ be the coefficient of $E$ in $D_W$,
and let $d$ be the multiplicity of $D'$ at the centre. The
exact divisor identity is
\[
 K_W+D_W=\sigma^*(K_{W'}+D')+(1+a-d)E.
                                                     \tag{C.IR.4}\label{cdp:eq:C.IR.4}
\]
Intersection with $E$ gives $a=d-1$, so \eqref{cdp:eq:C.IR.2} descends.
An exceptional curve in the vertical conductor is disjoint
from $H$. An exceptional curve outside the conductor has
$D_W\cdot E=1$, and therefore meets the horizontal part in
at most one simple point; it cannot identify two horizontal
points or join a horizontal component to the vertical
conductor. Thus the horizontal curves stay smooth, disjoint,
and disjoint from the remaining vertical conductor after
each blowdown. On the resulting minimal ruled bundle every
nonzero effective vertical divisor meets the horizontal
part. Its vertical conductor is consequently zero.

In the reverse blowup sequence a centre away from the current
horizontal divisor would have $d=0$ in \eqref{cdp:eq:C.IR.4} and create
coefficient $-1$ on the new exceptional curve. That coefficient
persists on its strict transform under later blowups, contrary
to effectiveness of the original $D_Z$. Every centre
therefore lies on the smooth coefficient-one horizontal
divisor and creates coefficient zero. Induction gives
\[
                             D_Z=H.               \tag{C.IR.5}\label{cdp:eq:C.IR.5}
\]
This argument did not require the map $h$ to factor through
any ruling blowdown.

We next prove that $\pi$ has no exceptional curves, including
horizontal ones. An exceptional curve $E$ of coefficient
zero in \eqref{cdp:eq:C.IR.5} satisfies $K_Z\cdot E\ge0$ on the minimal
resolution: adjunction and negative self-intersection give
this inequality, the sole contrary possibility being a
smooth rational $(-1)$-curve, which minimality excludes.
Equation \eqref{cdp:eq:C.IR.2} then forces $K_Z\cdot E=H\cdot E=0$.
Such a curve is disjoint from $H$. A positive-coefficient
exceptional curve is one of the mutually disjoint elliptic
curves $H_i$. Every connected exceptional fibre is therefore
disjoint from the strict conductor curves.

The normal conductor ideal on $T$ is the rank-one reflexive
sheaf $\omega_T\otimes\eta^*\omega_S^{-1}$, embedded in
$\mathcal O_T$ by finite duality. On an open set without
its conductor divisors it equals $\mathcal O_T$: equality
off a codimension-two set extends by reflexivity on a normal
surface. Thus $\eta$ is an isomorphism near the image of
every exceptional fibre just described. Such an image is
an isolated normal hypersurface singularity of $S$.
The isolated-singularity clause of \ref{cdp:st:A.5.17} excludes it from
this actual torus smoothing, including its full multiplicity
$m$ and its local unit. Hence $\pi$ is an isomorphism.
In particular $T=Z$ is smooth and its conductor is exactly
the reduced divisor $H$.

Let $G$ be the conductor image. Its ideal is the contraction
of the radical ideal of $H$ under the finite normalization,
so it is radical. Thus $G$ itself is reduced. Every local
point of $H$ has a single smooth branch, whose image is one
local branch of $G$. We claim that $G$ has only unibranch
germs and that its distinct global components are disjoint.
Otherwise choose two local branches of $G$ at a common point
and points of $T$ receiving the respective branches. These
points are different because $H$ is smooth at each point.
They correspond to different irreducible hypersurface
branches of $S$ in the smooth threefold. Two such branches
meet in a curve: their distinct equations generate an ideal
of height two in the regular three-dimensional local ring.
That intersection is a conductor-image branch received by
both points of $T$. Each of those points receives only one
branch, contradicting the chosen distinct images. This
proves the claim.

The generic transverse curve has conductor exponent one on
every normalized branch. The branchwise conductor formula
in the proof of \ref{cdp:st:C4e} shows that it has exactly two smooth
transverse branches: a singular branch has intrinsic
conductor exponent at least two, and for smooth branches
the exponent is the sum of their positive pairwise
intersection multiplicities. Consequently the finite map
$H\to G$ has total degree two over every irreducible component.

Fix such a component $G_j$, with normalization
$\nu:E_j\to G_j$. Each component of its inverse image is
elliptic. Riemann--Hurwitz gives $g(E_j)\le1$.
We exclude $E_j=\mathbf P^1$ by an actual local comparison.
In that case its inverse image must be one elliptic curve
of degree two; two degree-one maps from elliptic curves to
$\mathbf P^1$ cannot occur. Its four ramification points
have four distinct images on $\mathbf P^1$. \ref{cdp:st:A.5.14} gives
$p_a(G_j)\le2$, and all singularities of $G_j$ are unibranch.
There are therefore at most two normalization points above
its singular locus. At least one of the four branch points
lies over a smooth point of $G_j$.

At that point choose coordinates $(u,v)$ on $T$, with
$H=\{u=0\}$, so that the degree-two conductor map is
$z=v^2$. The normalization-conductor square is cartesian
as a square of local analytic rings. Because its conductor
is $(u)$, the surface local ring is exactly
\[
 \{a\in\mathbb C\{u,v\}:a(0,v)\in\mathbb C\{v^2\}\}
   =\mathbb C\{u,uv,v^2\}
   \simeq\mathbb C\{x,y,z\}/(x^2-y^2z).       \tag{C.IR.6}\label{cdp:eq:C.IR.6}
\]
Here $(x,y,z)=(uv,u,v^2)$, so every map in the comparison
is specified. The pinch germ in \eqref{cdp:eq:C.IR.6} is excluded by
\ref{cdp:st:A.5.15}. It follows that $E_j$ is elliptic.

Now \ref{cdp:st:A.5.14} bounds the total delta invariant of $G_j$ by one
and proves that its embedding in $Y$ is a local complete
intersection. On $E_j$ let $N$ be the dual of its pulled-back
conormal bundle, and let $\Xi$ be its normalization conductor
divisor. The determinant comparison from \ref{cdp:st:A.5.5} is
\[
 \det N=K_{E_j}(\Xi)\otimes\nu^*K_Y^{-1},
 \qquad \deg N=2\delta_j.                       \tag{C.IR.7}\label{cdp:eq:C.IR.7}
\]
The reduced equation of $S$ belongs to the ordinary square
of the ideal of $G_j$: the ideal is a regular sequence,
its powers are unmixed, and it has order at least two at
the generic point. Its quadratic jet is therefore a
holomorphic section of $\operatorname{Sym}^2N^*\otimes M$,
where $M=\nu^*\mathcal O_Y(S)$ has degree zero. Its determinant
is nonzero generically and has degree
\[
 \deg\bigl((\det N)^{-2}\otimes M^2\bigr)
                         =-4\delta_j.           \tag{C.IR.8}\label{cdp:eq:C.IR.8}
\]
Thus $\delta_j=0$. The curve $G_j$ is smooth elliptic and
the determinant, now a degree-zero holomorphic section,
is nowhere zero. The holomorphic Morse lemma with the
curve parameter gives the local reduced equation $xy=0$
at every point of $G_j$. This also proves that the
degree-two map $H\to G$ is etale everywhere.

There are no triple points, because $H$ is smooth and its
components are disjoint. Apply the already proved
normal-crossings comparison \ref{cdp:st:C.2.NC.1} to the actual full
fibre $mS$. It gives
\[
 0=e(Y_t)=m\,e(T\setminus H)=m\,e(T),            \tag{C.IR.9}\label{cdp:eq:C.IR.9}
\]
because $e(H)=0$. If $N_0$ is the number of ruling blowups
above a minimal bundle over the elliptic curve $B$, the
point-blowup formula gives $e(T)=N_0$. Hence $N_0=0$ and
$T$ is minimal. Finally normalization is an isomorphism
off $H$ and $G$, so
$e(S)=e(T)-e(H)+e(G)=0$. This proves every assertion,
with arbitrary full-fibre multiplicity retained.
\end{proof}

The elliptic ruled geometry in \ref{cdp:st:C.2.IRR.1} has a direct literature
comparison. Kamenova's Theorem 3.1 in \cite{KamenovaTori} treats
semistable Kahler degenerations with trivial total-space canonical
bundle; its elliptic ruled cases have elliptic double curves arranged
along a path or a cycle. Our finite normalization $T\to S$ specifies
the self-identification of the double curves in the irreducible case.
The proof above establishes this geometry from the actual conductor
and the full-fibre smoothing, retaining the multiplicity and proving
the necessary local pinch exclusion. We have not supplied a Kahler
form or a trivialization of the total-space canonical bundle needed
to apply that theorem directly to our tube. 


